\section*{Generative AI Usage Statement}
\textbf{[DRAFT --- authors: review, correct, and complete before submission; this paragraph was itself drafted by an AI assistant and records only what it did, not the authors' own tool use or verification process.]} An AI assistant (Claude, Anthropic) was used during manuscript preparation for: (i) checking the correctness of the paper's proofs, which surfaced several errors later corrected in the text (including a gap in the proof of the imperfect-observation proposition that affected its headline claim, and an unjustified periodicity claim in the recurring-audits theorem); (ii) drafting and revising prose throughout, including condensing the paper to meet the venue's page limit by moving full proofs and secondary robustness checks to the appendix; and (iii) writing and running the numerical scripts in \texttt{recurring\_sim/} that verify the corrected imperfect-observation result. All mathematical claims, numerical results, and text changes were reviewed before inclusion; authors should state here the specific extent and method of that review, and list any other generative-AI or routine spell/grammar tools used, which this statement does not cover.

\section*{Ethical Considerations Statement}
\paragraph{Ethical Considerations.} This paper is a theoretical and simulation-based contribution: it does not involve human subjects, does not collect new data, and does not audit or single out any specific firm. The empirical statistics we cite (Section~\ref{sec:calibration}) are aggregate, already-published figures from prior peer-reviewed field studies~\cite{wright2024null,gerchick2025auditing}; we do not use or reproduce any individual-level or re-identifiable data. Our central results are derived from a stylized formal model calibrated to illustrative parameter values, not fit to any specific deployed hiring or lending system, and should not be read as a factual claim about the compliance behavior of any named or identifiable organization.

\paragraph{Adverse Impacts.} Two results in this paper carry a real risk of being misapplied if lifted out of context. First, Proposition~\ref{prop:imperfectobs}'s finding that noisier observation can raise long-run quality in our model is not a recommendation to make real audits vaguer, less transparent, or less contestable; the model captures only one narrow welfare channel and abstracts away due-process and contestability values that a real audit regime must also serve, values that noisy or unaccountable audits can actively undermine. Second, the regulator-optimal threshold $\tau^\star = q^*(p_H)$ characterized in Section~\ref{sec:regulator} is a property of a stylized, calibrated model, not a ready-to-implement policy number; using it, or any single number from this paper, to directly set or defend a real regulatory threshold without independent empirical validation would overstate what the model supports. More generally, a paper that formalizes how audited firms can rationally reduce compliance investment could in principle be read by a firm as a playbook for exploiting these dynamics rather than as a diagnosis directed at regulators. We believe the diagnostic and regulatory-design value of naming this mechanism explicitly outweighs this risk, since the underlying incentive problem exists whether or not it is formalized here, but we flag the risk directly rather than leave it implicit.
